% Generated draft: requires mathematical and source review.
\documentclass[11pt]{article}
\usepackage[margin=25mm]{geometry}
\usepackage{fontspec}
\setmainfont{Latin Modern Roman}
\usepackage{amsmath,amssymb,mathtools,unicode-math}
\setmathfont{Latin Modern Math}
\usepackage{enumitem,needspace,xurl,hyperref}
\setlength{\parindent}{0pt}
\setlength{\parskip}{5pt}
\setlength{\emergencystretch}{3em}
\newcommand{\N}{\mathbb{N}}
\newcommand{\Z}{\mathbb{Z}}
\newcommand{\Q}{\mathbb{Q}}
\newcommand{\R}{\mathbb{R}}
\newcommand{\C}{\mathbb{C}}
\newcommand{\F}{\mathbb{F}}
\newcommand{\A}{\mathbb{A}}
\newcommand{\PP}{\mathbb{P}}
\newcommand{\E}{\mathbb{E}}
\newcommand{\eps}{\varepsilon}
\newcommand{\dd}{\,\mathrm{d}}
\begin{document}
\section*{UnsolvedMath ID: 30003492}
\textbf{Source problem code:} catalogue entry
\subsection{Definitions and mathematical statement}
A global field is a finite extension of \(\Q\), or a finite extension of \(\F_q(t)\), where \(\F_q\) is a finite field. Let \(G\) be a connected linear algebraic group over a global field \(k\). Write \(\mathbf G_m\) for the multiplicative group over \(k\), and write \(\#S\) for the cardinality of a set \(S\). The notation \(\tau(G)\) denotes the Tamagawa number, namely the volume of the adelic quotient \(G(k)\backslash G(\mathbb A_k)\) with its Tamagawa measure, subject to the normalization and convergence conventions needed for the class of groups under consideration; \(\mathbb A_k\) is the adele ring of \(k\). Those conventions are not specified in the supplied record.

The supplied question asks whether, for every global field \(k\) and every connected linear algebraic \(k\)-group \(G\), one has
\[
\tau(G)=\frac{\#\operatorname{Ext}^{1}(G,\mathbf G_m)}{\#X(G)}.
\]
The record does not define \(X(G)\) or specify the category and meaning of \(\operatorname{Ext}^{1}\). Thus this displayed question preserves the source formulation, but does not yet determine a fully specified mathematical assertion. In particular, the cardinalities in the quotient must be finite, and the denominator nonzero, for the displayed right-hand side to have its usual numerical meaning.

One possible interpretation, requiring confirmation from the cited source, is that \(\operatorname{Ext}^{1}(G,\mathbf G_m)\) is the group of equivalence classes of central extensions of algebraic \(k\)-groups
\[
1\longrightarrow\mathbf G_m\longrightarrow E\longrightarrow G\longrightarrow1,
\]
with equivalences inducing the identity on the kernel and quotient, and with Baer sum as group operation. Another possible intended convention is
\[
X(G)=\left\{\xi\in H^1(k,G):\xi|_{k_v}\text{ is the distinguished trivial class for every place }v\text{ of }k\right\},
\]
where \(k_v\) is the completion at \(v\), and \(H^1\) denotes the pointed set of isomorphism classes of torsors for the faithfully flat, finitely presented topology. For noncommutative \(G\), this is a pointed-set condition, not an assertion that \(H^1(k,G)\) is a group. These are explicitly provisional interpretations, not definitions established by the supplied record. The meanings of these two invariants, their finiteness, and the Tamagawa-measure conventions require human mathematical review before the universal identity can be treated as a precise statement.
\subsection{Short English statement}
The supplied problem asks whether a Tamagawa-number formula holds for every connected linear algebraic group over every global field. Its notation leaves \(X(G)\), the extension invariant, and the measure conventions insufficiently specified. The supplied status is partially solved; its dated literature assessment is recorded below without independent verification.
\subsection{Sources}
\begin{itemize}
\item Local UnsolvedMath registry, permanent ID \(30003492\), source problem code \texttt{catalogue entry}, title: \emph{Tamagawa Formula for Connected Linear Algebraic Groups}. The record gives proposed year 2017 and names no proposer. Its \texttt{corrected\_verified} statement label is accompanied by a description of notation, typography, and sentence-structure repairs; it does not establish mathematical correctness.
\item Supplied source citation: \emph{Algebraic Groups (2018)}, Oberwolfach Reports, DOI 10.4171/owr/2017/21, pp. 1281--1347. Supplied literature-source label: \emph{Oberwolfach Reports 2018}. \url{https://doi.org/10.4171/owr/2017/21}. The background describes this collection as arising from workshop open-problem discussions.
\item Supplied literature source: Rosengarten, \emph{Tamagawa numbers of pseudoreductive groups}. \url{https://msp.org/ant/2021/15-8/ant-v15-n8-p01-p.pdf}.
\item Supplied literature source: Rosengarten, \emph{arithmetic invariants of unipotent groups}. \url{https://arxiv.org/abs/1907.06225}.
\item Literature assessment dated exactly \texttt{2026-08-21}: the supplied curation labels the problem partially solved and states that the formula is proved for connected commutative groups and for all pseudoreductive groups in characteristic greater than three, while remaining small-characteristic exotic cases prevent a general solution. This is supplied dated curation derived from the Oberwolfach Reports record, not an independently verified assessment. The record explicitly calls for independent verification of time-sensitive status and cited claims.
\end{itemize}
\end{document}
